\documentclass[11pt,a4paper]{article}
\usepackage[T1]{fontenc}
\usepackage[utf8]{inputenc}
\usepackage[spanish,es-minimal]{babel}
\usepackage[margin=2.5cm]{geometry}
\usepackage{amsmath,booktabs,microtype}
\usepackage[hidelinks]{hyperref}
\hypersetup{pdftitle={El péndulo que se vacía},pdfsubject={Informe de laboratorio computacional}}

\title{El péndulo que se vacía\\[0.4em]\large Informe de laboratorio computacional}
\author{Proyecto \texttt{pendulum}}
\date{}

\begin{document}
\maketitle

\begin{abstract}
Se estudió por simulación un péndulo físico cuyo recipiente cilíndrico se
vacía por un orificio inferior. La altura del líquido se calculó con la ley
de Torricelli, y el movimiento angular se integró con Runge--Kutta de
cuarto orden. En el caso base, el vaciado duró \(96{,}33\,\mathrm{s}\) y el
período de pequeña amplitud alcanzó un máximo de \(1{,}6989\,\mathrm{s}\).
Una comparación con la aproximación adiabática, cinco comprobaciones del
modelo y un refinamiento del paso cuantificaron la consistencia de los
resultados numéricos.
\end{abstract}

\section{Objetivo}
Determinar cómo el vaciado modifica el período y la amplitud de un péndulo
físico. Implementar el mismo modelo en Python para reproducir los cálculos
y en JavaScript para explorar el efecto de los parámetros en el navegador.

\section{Modelo físico}
La botella oscila en un plano alrededor de un pivote fijo. La cuerda se
mantiene tensa y la botella se alinea con ella. Se aproxima el líquido
como un cilindro homogéneo de altura \(h(t)\) que gira con el recipiente.
El orificio se ubica sobre el eje del cilindro, a una distancia
\(s=L_c+H\) del pivote, donde \(L_c\) es la longitud de la cuerda y \(H\)
es la altura interior de la botella.

Con radio interior \(R\), diámetro del orificio \(d_a\), densidad \(\rho\)
y coeficiente de descarga \(C_d\), se definen \(A=\pi R^2\) y
\(a=\pi d_a^2/4\). La conservación de masa y la velocidad de descarga
dan \(A\dot h=-C_d a\sqrt{2gh}\). Para \(a>0\) y \(h_0>0\),
\[
 t_v=\frac{A}{C_da}\sqrt{\frac{2h_0}{g}},\qquad
 h(t)=h_0\max(0,1-t/t_v)^2.
\]
Para un orificio cerrado se mantiene \(h=h_0\). Al completarse el vaciado,
\(h=\dot h=0\). La masa de líquido es \(m_w=\rho Ah\).

Con \(d_b=L_c+H/2\) y masa de botella \(m_b\), el momento estático \(S\)
y el momento de inercia \(I\) respecto del pivote son
\begin{align*}
 S(h) &= m_b d_b + m_w(s-h/2),\\
 I(h) &= m_b(R^2/2+H^2/12+d_b^2)
       +m_w\big[R^2/4+h^2/12+(s-h/2)^2\big].
\end{align*}
El flujo de salida transporta momento angular a razón de
\(\dot m_{\mathrm{out}}s^2\dot\theta\), con
\(\dot m_{\mathrm{out}}=-\rho A\dot h\). El balance incluye ese flujo y el
torque gravitatorio \(-gS\sin\theta\):
\begin{align*}
 \frac{d}{dt}(I\dot\theta)+\dot m_{\mathrm{out}}s^2\dot\theta
      &= -gS\sin\theta,\\
 I\ddot\theta+\Gamma\dot\theta+gS\sin\theta &= 0,\qquad
 \Gamma=\rho A|\dot h|\big[h(2s-h)-R^2/4\big].
\end{align*}
Se obtiene la longitud equivalente \(L_{\mathrm{eq}}=I/S\) y el período
instantáneo de pequeña amplitud \(T_0=2\pi\sqrt{I/(gS)}\). El desarrollo
teórico del repositorio amplía las derivaciones y propone un protocolo
de contraste experimental \cite{teoria}.

\section{Implementación y procedimiento}
\begin{enumerate}
\item \texttt{fisica.py} calcula \(h\) y \(\dot h\) analíticamente y
evalúa \(S\), \(I\) y \(\Gamma\) en cada etapa de integración. Se resolvió
el sistema de primer orden \(\dot\theta=\omega\),
\(\dot\omega=-(\Gamma\omega+gS\sin\theta)/I\), con
\(\theta(0)=\theta_0\) y \(\omega(0)=0\). El paso fijo nominal de RK4 es
\(10^{-3}\,\mathrm{s}\).

\item Para estimar el período numérico, se interpoló linealmente el
instante de cada cruce ascendente de \(\theta=0\) y se restaron los
tiempos de dos cruces consecutivos. La amplitud de cada ciclo completo
se estimó con el mayor valor de \(|\theta|\) registrado en los pasos de
integración entre esos cruces.

\item \texttt{web/simulador.js} implementa las mismas ecuaciones y usa
Canvas para mostrar la oscilación y las series temporales. La velocidad
de animación cambia cuántos pasos se ejecutan por fotograma. La interfaz
permite variar las condiciones iniciales y comparar el período medido
por cruces con el período teórico.

\item Desde la raíz del repositorio se ejecutó
\texttt{python3 main.py --solo-verificar}. Ese comando ejecuta las cinco
comprobaciones e imprime los resultados del caso base. Para explorar el
efecto del vaciado en la interfaz, se puede ejecutar
\texttt{python3 main.py} y repetir el caso base con \(d_a=0\).
\end{enumerate}

\section{Resultados}
El caso base empleó \(H=h_0=0{,}250\,\mathrm{m}\),
\(R=0{,}040\,\mathrm{m}\), \(L_c=0{,}500\,\mathrm{m}\),
\(m_b=0{,}040\,\mathrm{kg}\), \(d_a=5\,\mathrm{mm}\), \(C_d=0{,}60\),
\(\theta_0=12^\circ\), \(\rho=1000\,\mathrm{kg/m^3}\) y
\(g=9{,}81\,\mathrm{m/s^2}\).

\begin{table}[htbp]
\centering
\caption{Predicciones del modelo para el caso base. Los períodos de esta
tabla se calculan con la aproximación de pequeña amplitud.}
\begin{tabular}{lc}
\toprule
Magnitud & Valor \\
\midrule
Tiempo de vaciado \(t_v\) & \(96{,}33\,\mathrm{s}\) \\
Período inicial & \(1{,}5973\,\mathrm{s}\) \\
Período máximo & \(1{,}6989\,\mathrm{s}\) \\
Instante y altura en el máximo & \(61{,}1\,\mathrm{s}\), \(33{,}4\,\mathrm{mm}\) \\
Período con botella vacía & \(1{,}5981\,\mathrm{s}\) \\
Factor adiabático final \(\Theta_f/\Theta_0\) & \(0{,}71984\) \\
\bottomrule
\end{tabular}
\end{table}

\begin{table}[htbp]
\centering
\caption{Comprobaciones ejecutadas con el modelo de referencia en Python.}
\small
\begin{tabular}{@{}p{0.39\linewidth}p{0.36\linewidth}p{0.17\linewidth}@{}}
\toprule
Comprobación & Resultado & Tolerancia \\
\midrule
Período de amplitud finita con orificio cerrado &
\(1{,}601694\,\mathrm{s}\), error \(2{,}0\times10^{-8}\) &
\(2\times10^{-4}\) \\
Derivadas de \(S\) e \(I\) &
Error relativo máximo \(6{,}3\times10^{-12}\) & \(10^{-6}\) \\
Amplitud para \(d_a=8\,\mathrm{mm}\) &
RK4 \(0{,}72443\), WKB \(0{,}71984\), error \(0{,}64\,\%\) & \(1\,\%\) \\
Amplitud para \(d_a=6\) y \(12\,\mathrm{mm}\) &
Factores \(0{,}72043\) y \(0{,}72875\), diferencia \(1{,}14\,\%\) & \(2\,\%\) \\
Límite de masa puntual &
Variaciones del período y la amplitud inferiores a \(0{,}001\,\%\) &
\(0{,}5\,\%\) \\
\bottomrule
\end{tabular}
\end{table}

Para evaluar la sensibilidad al paso se ejecutó \texttt{fisica.integrar}
con los parámetros base, \texttt{t\_final=100} y \texttt{muestreo=1},
variando \texttt{dt} como se indica en la tabla~\ref{tab:paso}.
El ángulo se tomó de la última muestra, \texttt{r['theta'][-1]}.

\begin{table}[htbp]
\centering
\caption{Sensibilidad al paso de integración en el caso base.}
\label{tab:paso}
\begin{tabular}{cc}
\toprule
Paso \(\Delta t\) (\(\mathrm{ms}\)) & \(\theta(100\,\mathrm{s})\) (\(\mathrm{rad}\)) \\
\midrule
\(2\)   & \(0{,}055684244362\) \\
\(1\)   & \(0{,}055684243009\) \\
\(0{,}5\) & \(0{,}055684242924\) \\
\bottomrule
\end{tabular}
\end{table}

\section{Discusión}
El máximo del período supera en un \(6{,}36\,\%\) el valor inicial.
Surge de la evolución del cociente \(I/S\): retirar agua cambia tanto
la distribución de masa que determina \(I\) como el brazo gravitatorio
que determina \(S\). Al quedar vacía la botella, el período se aproxima de
nuevo a su valor inicial.

La aproximación de Wentzel--Kramers--Brillouin (WKB) predice
\[
 \frac{\Theta_f}{\Theta_0}=\sqrt{\frac{\Omega_0}{\Omega_f}}\,
 \exp\!\left[-\frac12\int_0^{t_v}\frac{\Gamma}{I}\,dt\right],
 \qquad \Omega=\sqrt{gS/I}.
\] Al sustituir \(dt=dh/\dot h\), la integral depende
de la geometría y del llenado inicial. La diferencia del \(1{,}14\,\%\)
entre orificios cuantifica la variación de las estimaciones numéricas
para dos velocidades de vaciado. Se usó el último ciclo completo
disponible antes de \(0{,}995\,t_v\) para esas estimaciones.

El refinamiento de \(1\) a \(0{,}5\,\mathrm{ms}\) modificó el ángulo a
\(100\,\mathrm{s}\) en aproximadamente \(8{,}5\times10^{-11}\,\mathrm{rad}\).
Esa comparación mide la sensibilidad numérica en el caso base. Los
resultados proceden de una simulación con descarga cuasiestacionaria,
sin oleaje interno, fricción del pivote ni arrastre del aire. Su
contraste físico requeriría medir \(h(t)\) y \(\theta(t)\) en un montaje
instrumentado.

\section{Conclusiones}
En el caso base simulado, el período presenta un máximo antes de que la
botella se vacíe y la aproximación adiabática predice una amplitud final
cercana al \(72\,\%\) de la inicial. Las cinco comprobaciones pasaron
dentro de las tolerancias definidas, y el refinamiento del paso confirmó
la estabilidad del ángulo final para los pasos examinados.

\begin{thebibliography}{9}
\bibitem{teoria}
Proyecto \texttt{pendulum}, \emph{El péndulo que se vacía},
desarrollo teórico incluido en \path{documento/pendulo_vaciado.tex}.
\end{thebibliography}
\end{document}
